\uuid{MLC3}
\titre{ Effet du taux d'apprentissage sur la descente de gradient }
\niveau{L3}
\module{Analyse de données}
\chapitre{Réseaux de neurones}
\sousChapitre{Descente de gradient, hyperparamètres}
\theme{Learning rate, divergence, oscillation}
\auteur{Maxime Nguyen}
\datecreate{2026-03-19}
\organisation{AMSCC}
\difficulte{2}

\contenu{
\texte{
  On se place dans la situation suivante : $w = 3$, gradient $= 4$.
}
\begin{enumerate}

\item
  \question{Avec $\alpha = 0{,}1$ : quelle est la nouvelle valeur de $w$ ?}
  \reponse{$w \leftarrow 3 - 0{,}1 \times 4 = 2{,}6$.}

\item
  \question{Avec $\alpha = 1$ : quelle est la nouvelle valeur de $w$ ?}
  \reponse{$w \leftarrow 3 - 1 \times 4 = -1$.}

\item
  \question{Avec $\alpha = 10$ : quelle est la nouvelle valeur de $w$ ? Que se passe-t-il si le gradient reste $\approx 4$ à l'itération suivante ?}
  \reponse{
    $w \leftarrow 3 - 10 \times 4 = -37$.
    Si le gradient reste $\approx 4$, la mise à jour suivante donne $w \leftarrow -37 - 40 = -77$.
    Le paramètre $w$ diverge : le coût augmente au lieu de diminuer.
    Un taux d'apprentissage trop élevé empêche la convergence.
  }

\end{enumerate}
}
